\documentclass[11pt]{article}
\usepackage[margin=1in]{geometry}
\usepackage{amsmath,amssymb,amsthm}
\usepackage{booktabs}
\usepackage{hyperref}

\newtheorem{theorem}{Theorem}
\newtheorem{lemma}[theorem]{Lemma}
\newtheorem{definition}[theorem]{Definition}
\newtheorem{remark}[theorem]{Remark}

\title{Multi-Threshold Pass-Count Preservation under Bounded Score Perturbations}
\author{Research Architecture Runtime \\ \texttt{operator/multi-threshold}}
\date{2026-07-25}

\begin{document}
\maketitle

\begin{abstract}
We study the multi-threshold pass-count
\(C_{\mathbf{T}}(x)=\lvert\{i:x\ge T_i\}\rvert\) for a finite threshold list
\(\mathbf{T}\) under \(\lvert x'-x\rvert\le\varepsilon\).
The count is invariant whenever every coordinate is outside its half-open
unstable band inherited from scalar thresholding.
Sharpness follows by flipping any single unstable cut.
Lean formalization is Mathlib $\mathbb{R}$ (\texttt{REAL\_MATHLIB})
(\texttt{LEAN\_FULL}; \texttt{REAL\_MATHLIB}).
\end{abstract}

\section{Introduction}
Multi-threshold selection aggregates many AboveThreshold bits into a Nat count.
It is the natural finite extension of scalar thresholding and a reduction target
for ordered cut-point / bucket operators.

\section{Operator definition}
\begin{definition}
For finite \(x\in\mathbb{R}\) and \(\mathbf{T}=(T_0,\ldots,T_{n-1})\),
\[
C_{\mathbf{T}}(x)=\sum_{i=0}^{n-1}\mathbf{1}\{x\ge T_i\}.
\]
Equality passes on each coordinate.
\end{definition}

\section{Perturbation model}
Admissible neighbors satisfy \(\lvert x'-x\rvert\le\varepsilon\) with \(\varepsilon\ge 0\).

\section{Instability}
Coordinate \(i\) is unstable when \(x\in[T_i-\varepsilon,T_i+\varepsilon)\).
The multi-threshold count is unstable whenever any coordinate is.

\section{Main theorems}
\begin{theorem}[Multi-threshold preservation]\label{thm:main}
If \(\lvert x'-x\rvert\le\varepsilon\) and for every \(i\) either
\(x\ge T_i+\varepsilon\) or \(x<T_i-\varepsilon\), then
\(C_{\mathbf{T}}(x')=C_{\mathbf{T}}(x)\).
\end{theorem}

\begin{theorem}[Sharpness]\label{thm:sharp}
If some \(j\) satisfies \(x\in[T_j-\varepsilon,T_j+\varepsilon)\), then there exists
admissible \(x'\) with \(C_{\mathbf{T}}(x')\ne C_{\mathbf{T}}(x)\)
(equivalently: the singleton list \([T_j]\) flips).
\end{theorem}

\section{Proofs}
\begin{proof}[Proof of Theorem~\ref{thm:main}]
Induct on the threshold list. The empty list has count \(0\).
For \(T::\mathbf{T}'\), the head bit is preserved by scalar threshold preservation;
the tail count is preserved by the inductive hypothesis on the remaining stable
coordinates.
\end{proof}

\begin{proof}[Proof of Theorem~\ref{thm:sharp}]
On the singleton \([T]\), push \(x\) by \(\pm\varepsilon\) toward the cut.
If \(x\ge T\) then \(x'=x-\varepsilon\) fails; if \(x<T\) then \(x'=x+\varepsilon\) passes.
\end{proof}

\section{Comparison with prior operators}
Relative to Threshold: same equality convention and half-open unstable band;
new selected object is a Nat pass-count; proof pattern is coordinatewise reduction.
Relative to Argmax (reference): different geometry (cuts vs margin among alternatives).

\section{Lean formalization summary}
Module \texttt{Research.Operators.MultiThreshold.Preservation}:
\texttt{multi\_threshold\_preservation}, \texttt{multi\_threshold\_sharpness}.
Domain: Mathlib $\mathbb{R}$ (\texttt{REAL\_MATHLIB}).

\section{Certificate}
\texttt{LEAN\_FULL} at
\texttt{lean/certificates/multi-threshold/multi-threshold-count-preservation/}.

\section{Limitations}
\begin{itemize}
\item Mathlib \(\mathbb{R}\) (\texttt{REAL\_MATHLIB}).
\item Unordered pass-count; ordered bucket indices are a separate operator.
\item No DP / privacy claims.
\end{itemize}

\section*{References}
\begin{thebibliography}{9}
\bibitem{repo}
Repository: \texttt{implementation/src/operators/multi\_threshold/}.
\end{thebibliography}

\end{document}
